\chapter{Panic Attacks And Panic Disorder}

\medskip
\noindent\textit{\textbf{Under review.} This chapter is an AI-assisted educational draft; it is not a substitute for professional medical advice.}

\section*{Definition and Overview}
Panic Attacks And Panic Disorder is a mental-health topic involving thoughts, emotions, behaviour, relationships, or functioning. A diagnosis requires assessment of duration, distress, impairment, safety, and possible physical or substance-related causes.

\section*{Clinical Features}
Possible features include changes in mood, anxiety, sleep, energy, concentration, behaviour, perception, relationships, or ability to function.

\section*{When to Seek Medical Care}
Seek medical review for new, persistent, worsening, or function-limiting symptoms. Call 000 for severe breathing difficulty, collapse, sudden neurological change, severe chest or abdominal pain, or any rapidly worsening illness.

\section*{Diagnosis and Investigations}
Assessment is clinical and may include risk assessment, mental-state examination, physical review, medication and substance history, and targeted tests to exclude medical causes.

\section*{Management}
Treatment may include evidence-based psychological therapy, medicines when indicated, social support, safety planning, and coordinated GP or specialist care.

\section*{Complications}
Complications depend on the underlying cause and may include progression, effects on other organs, treatment adverse effects, or reduced quality of life. Early assessment can reduce preventable harm.

\section*{Prognosis and Self-Care}
Outcome depends on the cause, severity, complications, and response to treatment. Follow the care plan, keep appointments, avoid known triggers, and seek help if symptoms worsen.

\section*{Expanded clinical context}
Panic attacks and panic disorder is a mental-health, behavioural, or support topic. Interpretation should account for age, baseline health, medicines, comorbidities, goals, and access to care. This chapter supports discussion with a qualified clinician and does not establish a diagnosis.

\section*{Assessment and decision points}
Review onset, duration, triggers, sleep, substance use, physical health, developmental context, social stressors, and immediate safety. Ask about self-harm or harm from others when appropriate. Document the main concern, time course, functional impact, relevant risks, and red flags. Use targeted investigations when their result is likely to change management.

\section*{Management and follow-up}
Care may include education, psychological therapy, practical support, treatment of coexisting conditions, medication when indicated, and a shared follow-up plan. Explain expected improvement, when reassessment is needed, and how follow-up can be accessed. Avoid starting, stopping, or sharing prescription medicines without professional advice.

\section*{Safety-netting}
Seek urgent help for immediate danger, suicidal intent, psychosis, severe confusion, inability to meet basic needs, or rapidly worsening symptoms.

\section*{Practical prevention}
Use plain-language information, supportive relationships, regular routines, and professional review when symptoms persist or impair daily life. Keep written instructions and seek review if the topic recurs, changes, or does not improve as expected.

\section*{Limitations}
This educational expansion should be checked against current local guidance and authoritative topic-specific sources.
\clearpage
